\documentclass[11pt]{article}
\usepackage[T1]{fontenc}
\usepackage[a4paper,margin=1in]{geometry}
\usepackage{amsmath,amssymb,amsthm}
\usepackage{graphicx}
\usepackage[hidelinks]{hyperref}
\usepackage{microtype}

\newtheorem{theorem}{Theorem}
\newtheorem{lemma}{Lemma}
\newtheorem{remark}{Remark}

\newcommand{\R}{\mathbb R}
\newcommand{\E}{\mathcal E}
\newcommand{\Q}{\mathcal Q}

\setlength{\parindent}{0pt}
\setlength{\parskip}{0.58em}

\title{The Exact Constant in a Weighted Critical Sobolev Inequality}
\author{Run \texttt{fa\_banach\_001} --- lane 14}
\date{9 August 2026}

\begin{document}
\maketitle

\begin{center}
\fbox{\parbox{0.91\textwidth}{
\textbf{Candidate full result; likely valid.}
For every $\alpha>0$, the best constant asked for in Remark~3 of
arXiv:2412.15490 is independent of $\alpha$ and equals the sharp Euclidean
three-dimensional constant
\[
 C_{2\alpha,6}=\sqrt3\left(\frac{\pi}{2}\right)^{2/3}.
\]
The key is an exact transformation to the ordinary Sobolev quotient on a
flat angle-excess cone, followed by sharp three-dimensional
Cartan--Hadamard isoperimetry.
}}
\end{center}

\section*{Source question}

Giang, Tri, and Tuan~\cite{source} define
\[
 C_{2\alpha,6}
 =\inf_{u\ne0}
 \frac{\left(\int_{\R^3}(|\nabla_xu|^2+|x|^{2\alpha}|u_y|^2)
 \,dx_1dx_2dy\right)^{1/2}}
 {\left(\int_{\R^3}|x|^{2\alpha}|u|^6\,dx_1dx_2dy\right)^{1/6}},
\]
where the infimum is over their space $W^{1,2\alpha,6}_0(\R^3)$, the
completion of $C_c^\infty(\R^3)$ in the sum of these two norms.  Remark~3,
on source PDF page~10, asks for its exact value.

\begin{figure}[ht]
\centering
\includegraphics[width=\textwidth]{figures/open_problem_crop.png}
\caption{Remark~3 in the source paper, including the surrounding
normalization.}
\end{figure}

\textbf{Transcription.} ``Our proof only gives a lower bound for the best
constant $C_{2\alpha,6}$.  The exact value of $C_{2\alpha,6}$ remains an open
question.''

\section*{Proof intuition}

The exponent and weights look anisotropic in $(x_1,x_2,y)$, but they are
exactly the Jacobian and metric coefficients of a flat cone.  If
$r=|x|$, stretch the polar angle and replace $r$ by
$r^{\alpha+1}/(\alpha+1)$.  Both integrals in the quotient then become the
ordinary unweighted $D^{1,2}\to L^6$ integrals on
\[
 X_\alpha=C_{2\pi(\alpha+1)}\times\R.
\]
The total cone angle exceeds $2\pi$, so the singular axis has angle excess,
the metric manifestation of nonpositive curvature.

Kleiner's theorem applies to smooth Cartan--Hadamard three-manifolds, whereas
$X_\alpha$ has a singular axis.  Two elementary devices bridge this gap.
First, a logarithmic cutoff removes the axis with vanishing energy.  Second,
every compact subset of the remaining cone exterior sits isometrically in a
smooth complete nonpositively curved surface of revolution times $\R$.
Kleiner's sharp isoperimetric inequality therefore gives the Euclidean sharp
Sobolev lower bound on the cone.  The reverse inequality follows because
arbitrarily large genuinely Euclidean balls can be placed far from the axis.

\section*{Statement}

Let
\[
 C_{\mathrm E,3}
 :=\inf_{0\ne w\in C_c^\infty(\R^3)}
 \frac{\|\nabla w\|_{L^2(\R^3)}}{\|w\|_{L^6(\R^3)}}
 =\sqrt3\left(\frac{\pi}{2}\right)^{2/3}.
\]

\begin{theorem}\label{thm:main}
For every $\alpha>0$,
\[
 C_{2\alpha,6}=C_{\mathrm E,3}
 =\sqrt3\left(\frac{\pi}{2}\right)^{2/3}.
\]
Equivalently, every $u\in W^{1,2\alpha,6}_0(\R^3)$ satisfies
\[
 C_{\mathrm E,3}
 \left(\int_{\R^3}|x|^{2\alpha}|u|^6\right)^{1/6}
 \le
 \left(\int_{\R^3}
 (|\nabla_xu|^2+|x|^{2\alpha}|u_y|^2)\right)^{1/2},
\]
and no larger constant works.
\end{theorem}

\section*{Exact transformation to a flat cone}

Put $a=\alpha+1>1$.  Let $C_{2\pi a}$ be the flat two-dimensional cone
whose polar coordinates are $(\rho,\varphi)$ with
$\rho>0$ and $\varphi\in\R/(2\pi a\mathbb Z)$, and set
\[
 X_a=C_{2\pi a}\times\R,
 \qquad
 g_a=d\rho^2+\rho^2d\varphi^2+dz^2.
\]
The singular axis is $A=\{\rho=0\}\times\R$.

On $\R^3\setminus\{x=0\}$ use cylindrical coordinates $(r,\theta,y)$ and
define
\begin{equation}\label{eq:T}
 \rho=\frac{r^a}{a},\qquad \varphi=a\theta,\qquad z=y.
\end{equation}
Because $\theta$ is taken modulo $2\pi$, $\varphi$ is taken modulo $2\pi a$;
thus \eqref{eq:T} is a diffeomorphism onto $X_a\setminus A$.
For $v(\rho,\varphi,z)=u(r,\theta,y)$,
\[
 d\rho=r^\alpha dr,\qquad d\varphi=a\,d\theta,
\]
and consequently
\begin{equation}\label{eq:measure}
 \rho\,d\rho\,d\varphi\,dz
 =r^{2\alpha+1}\,dr\,d\theta\,dy
 =|x|^{2\alpha}\,dx_1dx_2dy.
\end{equation}
Moreover
\[
 u_r=r^\alpha v_\rho,\qquad u_\theta=a v_\varphi,\qquad u_y=v_z.
\]
Using $dx_1dx_2dy=r\,dr\,d\theta\,dy$ and
$r^{2a}=a^2\rho^2$ gives
\begin{align}
 &\int_{\R^3}
 \left(u_r^2+\frac{u_\theta^2}{r^2}+r^{2\alpha}u_y^2\right)
 r\,dr\,d\theta\,dy \notag\\
 &\qquad=
 \int_{X_a}\left(v_\rho^2+\frac{v_\varphi^2}{\rho^2}+v_z^2\right)
 \rho\,d\rho\,d\varphi\,dz
 =\int_{X_a}|\nabla_{g_a}v|^2\,dV_{g_a}.             \label{eq:energy}
\end{align}
Equations \eqref{eq:measure} and \eqref{eq:energy} show that the source
quotient is exactly the ordinary critical Sobolev quotient on $X_a$.

\section*{The sharp lower bound on the cone}

We first record the smooth-completion device.

\begin{lemma}[Smooth Cartan--Hadamard completion]\label{lem:smooth}
If $K\subset X_a\setminus A$ is compact, a neighborhood of $K$ embeds
isometrically into a complete simply connected smooth three-manifold of
nonpositive sectional curvature.
\end{lemma}

\begin{proof}
Write $\varphi=a\vartheta$, where $\vartheta$ is modulo $2\pi$.  Then
\[
 g_a=d\rho^2+a^2\rho^2d\vartheta^2+dz^2.
\]
Choose a smooth nondecreasing function $\psi:[0,\infty)\to[1,a]$ which is
$1$ near $0$ and equals $a$ after a short transition interval $[0,s_0]$.
Set
\[
 h(s)=\int_0^s\psi(t)\,dt.
\]
The metric
\[
 ds^2+h(s)^2d\vartheta^2+dz^2
\]
is smooth at $s=0$, complete, and defined on a manifold diffeomorphic to
$\R^3$.  Its only possibly nonzero sectional curvature is the curvature of
the surface-of-revolution factor,
\[
 -\frac{h''(s)}{h(s)}=-\frac{\psi'(s)}{h(s)}\le0.
\]
For $s\ge s_0$ one has $h(s)=as+b=a(s-c)$ for a constant $c\ge0$.
Hence this exterior, after setting $\rho=s-c$, is exactly the cone exterior
with metric $d\rho^2+a^2\rho^2d\vartheta^2+dz^2$.

By shortening the transition interval, its matching cone radius
$h(s_0)/a$ can be made arbitrarily small.  Choose it below
$\min_K\rho>0$.  The resulting exterior contains an isometric neighborhood
of $K$.
\end{proof}

Kleiner's three-dimensional Cartan--Hadamard theorem~\cite{kleiner} states
that every bounded smooth domain in a complete simply connected smooth
three-manifold of nonpositive sectional curvature satisfies the sharp
Euclidean isoperimetric inequality.  The standard coarea and Schwarz
rearrangement argument then yields, for every compactly supported smooth
function $w$ on such a manifold,
\begin{equation}\label{eq:CHsob}
 C_{\mathrm E,3}\|w\|_6\le\|\nabla w\|_2.
\end{equation}
For completeness, the implication is: replace every superlevel set of
$|w|$ by a Euclidean ball of equal volume; Kleiner's perimeter inequality
and coarea decrease the Dirichlet integral, while the $L^6$ norm is
preserved, and the Euclidean Aubin--Talenti inequality finishes the estimate.
This sharp implication is also recorded explicitly in
Muratori--Soave~\cite{muratori-soave}.

Let $v\in C_c^\infty(X_a\setminus A)$.  Apply
Lemma~\ref{lem:smooth} to its support and then \eqref{eq:CHsob}; because the
embedding is isometric on a neighborhood of the support,
\begin{equation}\label{eq:away}
 C_{\mathrm E,3}\|v\|_{L^6(X_a)}
 \le\|\nabla v\|_{L^2(X_a)}.
\end{equation}

It remains to remove the restriction away from $A$.  For
$0<\varepsilon<1/2$, choose a radial cutoff $\chi_\varepsilon$ which is zero
for $\rho\le\varepsilon^2$, one for $\rho\ge\varepsilon$, and linear in
$\log\rho$ between these radii.  Thus
\[
 |\chi_\varepsilon'(\rho)|
 \le\frac{1}{\rho\log(1/\varepsilon)}
 \quad(\varepsilon^2<\rho<\varepsilon).
\]
On every bounded $z$-interval,
\begin{equation}\label{eq:cut}
 \int |\nabla\chi_\varepsilon|^2\,dV_{g_a}
 \le \frac{C}{\log^2(1/\varepsilon)}
 \int_{\varepsilon^2}^{\varepsilon}\frac{d\rho}{\rho}
 =\frac{C}{\log(1/\varepsilon)}\longrightarrow0.
\end{equation}
If $v$ is bounded, smooth away from the axis, and compactly supported, then
$v_\varepsilon=\chi_\varepsilon v$ satisfies
\[
 v_\varepsilon\to v\text{ in }L^6(X_a),
 \qquad
 \nabla v_\varepsilon\to\nabla v\text{ in }L^2(X_a),
\]
by \eqref{eq:cut}, dominated convergence, and the product rule.
The transform of every $u\in C_c^\infty(\R^3)$ has these properties: it is
bounded at the axis and smooth off it.  Passing to the limit in
\eqref{eq:away}, then using \eqref{eq:measure}--\eqref{eq:energy}, proves
\[
 C_{2\alpha,6}\ge C_{\mathrm E,3}.
\]

\section*{Matching upper bound and the value}

The flat cone $X_a$ contains isometric copies of Euclidean balls of
arbitrarily large radius: for a prescribed radius $R$, place the center at
cone distance much larger than $R$ from the axis and choose an angular chart
which does not meet the identified seam.  Therefore every compactly
supported Euclidean test function can, after a harmless dilation and
translation, be inserted into $X_a$ with its Sobolev quotient unchanged.
Taking a sequence of truncated Aubin--Talenti profiles whose Euclidean
quotients tend to $C_{\mathrm E,3}$ gives
\[
 C_{2\alpha,6}\le C_{\mathrm E,3}.
\]
Together with the lower bound this proves equality.

Finally, for the Euclidean extremal $U(x)=(1+|x|^2)^{-1/2}$,
\[
 \int_{\R^3}U^6\,dx=\frac{\pi^2}{4},
 \qquad
 \int_{\R^3}|\nabla U|^2\,dx=\frac{3\pi^2}{4}.
\]
Hence
\[
 C_{\mathrm E,3}
 =\frac{(3\pi^2/4)^{1/2}}{(\pi^2/4)^{1/6}}
 =\sqrt3\left(\frac{\pi}{2}\right)^{2/3},
\]
which completes the proof of Theorem~\ref{thm:main}.  Density extends the
inequality from $C_c^\infty$ to the source space
$W^{1,2\alpha,6}_0(\R^3)$.

\section*{Scope, novelty, and review status}

The theorem determines exactly the constant asked for in Remark~3.  It does
not classify extremizing sequences or assert that the infimum is attained.

A bounded search through 9 August 2026 used arXiv:2412.15490, the exact
paper title and open-question phrase, $C_{2\alpha,6}$, `weighted Sobolev best
constant', `Grushin sharp Sobolev constant', `flat cone Sobolev constant',
and the source authors.  It found the source and general sharp
Cartan--Hadamard Sobolev literature, but no later exact answer or this cone
identification.  This is not a comprehensive priority search.

The current verdict is \textbf{likely valid}; human review is recommended.
The most important checks are the exact energy identity, the smooth
completion lemma, and the zero-capacity cutoff at the cone axis.

\begin{thebibliography}{9}
\bibitem{source}
T. H. Giang, N. M. Tri, and D. A. Tuan,
\emph{On some Sobolev and P\'olya--Szeg\"o type inequalities with weights
and applications}, arXiv:2412.15490.

\bibitem{kleiner}
B. Kleiner,
\emph{An isoperimetric comparison theorem},
Invent. Math. \textbf{108} (1992), 37--47.
\href{https://doi.org/10.1007/BF02100598}{doi:10.1007/BF02100598}.

\bibitem{muratori-soave}
M. Muratori and N. Soave,
\emph{Some rigidity results for Sobolev inequalities and related PDEs on
Cartan--Hadamard manifolds}, Ann. Sc. Norm. Super. Pisa Cl. Sci. (5)
\textbf{24} (2023), 59--103; arXiv:2103.08240.

\bibitem{talenti}
G. Talenti,
\emph{Best constant in Sobolev inequality},
Ann. Mat. Pura Appl. (4) \textbf{110} (1976), 353--372.
\end{thebibliography}

\end{document}
